% texte.tex — f176v, vv. 1169–1210 (conversion initiale des travaux antérieurs)
\vers[n=1169]{}{}{ἀπὸ τῆς στρατιᾶς παροῦσιν ὑμῶν διαβαλῶ.}{à vos maris, à leur retour de l'armée.}
\vers[n=1170]{ΚΟ.}{Coryphée}{τὰ μὲν παρ᾽ ἡμῶν ἴσθι σοι πεπεισμένα·}{Pour ce qui est de nous, sache que nous acquiesçons ;}
\vers[n=1171]{}{}{τὸν βάρβαρον δὲ τοῦτον αὐτὸς πεῖθε σύ.}{mais ce barbare, c'est à toi de le convaincre.}
\vers[n=1172]{ΕΥ.}{Euripide}{ἐμὸν ἔργον ἐστίν· καὶ σὸν ὦλάφιον ἅ σοι}{Ça, c'est mon affaire. \didasc{À la danseuse} La tienne, Élaphion, est de songer à t'y prendre comme je t'ai dit chemin faisant.}
\vers[n=1173]{}{}{καθ᾽ ὁδὸν ἔφραζον ταῦτα μεμνῆσθαι ποιεῖν.}{Ainsi donc, passe d'abord à l'autre bout, et reviens en courant sur tes pointes.}
\vers[n=1174]{}{}{πρῶτον μὲν οὖν δίελθε κἀνακάλπασον.}{\didasc{Au flûtiste :} Toi, Térédon, joue-nous sur la flûte une Persique.}
\vers[n=1175]{}{}{σὺ δ᾽ ὦ Τερηδὸν ἐπαναφύσα Περσικόν.}{\didasc{Térédon joue de la flûte, pendant que la jeune fille danse.}}
\vers[n=1176]{ΤΟ.}{Archer}{τί τὸ βόμβο τοῦτο; κῶμό τις ἀνεγεῖρί μοι;}{\didasc{Se réveillant.} Qu'est-ce que ce ronflement? c'est un comos qui me réveille ?}
\vers[n=1177]{ΕΥ.}{Euripide}{ἡ παῖς ἔμελλε προμελετᾶν ὦ τοξότα.}{Cette enfant, archer, pensait s'essayer.}
\vers[n=1178]{}{}{ὀρχησομένη γὰρ ἔρχεθ᾽ ὡς ἄνδρας τινάς.}{Elle s'en va danser chez des particuliers.}
\vers[n=1179]{ΤΟ.}{Archer}{ὀρκῆσι καὶ μελετῆσι, οὐ κωλῦσ᾽ ἐγώ.}{Qu'elle danse, qu'elle s'essaie, je ne m'y oppose pas.}
\vers[n=1180]{}{}{ὠς ἐλαπρός, ὤσπερ ψύλλο κατὰ τὸ κῴδιο.}{Qu'elle est vive! On dirait une puce sur une peau de bique.}
\vers[n=1181]{ΕΥ.}{Euripide}{φέρε θοἰμάτιον ἄνωθεν ὦ τέκνον τοδί·}{Tends-moi ce vêtement, ma fille, en le tirant par en haut.}
\vers[n=1182]{}{}{καθιζομένη δ᾽ ἐπὶ τοῖσι γόνασι τοῦ Σκύθου}{Assieds-toi sur les genoux du Scythe,}
\vers[n=1183]{}{}{τὼ πόδε πρότεινον, ἵν᾽ ὑπολύσω.}{et avance les pieds, que je les déchausse.}
\vers[suite]{ΤΟ.}{Archer}{ναῖκι ναὶ}{Si fait, si,}
\vers[n=1184]{}{}{κάτησο κάτησο, ναίκι ναί, τυγάτριον.}{mets-toi, si fait, ma fillette.}
\vers[n=1185]{}{}{οἴμ᾽ ὠς στέριπο τὸ τιττί᾽, ὤσπερ γογγυλί.}{\didasc{Pendant qu'Euripide ôte les chaussures à la danseuse.} Ouais! est-il ferme, ce téton ! une vraie rave. \didasc{La jeune fille déchaussée et nue danse au son de la flûte.}}
\vers[n=1186]{ΕΥ.}{Euripide}{αὔλει σὺ θᾶττον· ἔτι δέδοικας τὸν Σκύθην;}{Joue donc plus vite. As-tu encore peur du Scythe ?}
\vers[n=1187]{ΤΟ.}{Archer}{καλό γε τὸ πυγή. κλαῦσί γ᾽ ἢν μὴ 'νδον μένῃς.}{Ah! que voilà de belles fesses ! — Gare à toi si tu ne restes dans ton coin.}
\vers[n=1188]{}{}{εἶεν· καλὴ τὸ σκῆμα περὶ τὸ πόστιον.}{\didasc{Il s'arrange.} — À la bonne heure ! Elle a repris de la tenue, la broquette.}
\vers[n=1189]{ΕΥ.}{Euripide}{καλῶς ἔχει. λαβὲ θοἰμάτιον· ὥρα 'στὶ νῷν}{\didasc{À la danseuse}. C'est assez. Reprends ton vêtement. Il est temps que}
\vers[n=1190]{}{}{ἤδη βαδίζειν.}{nous partions.}
\vers[suite]{ΤΟ.}{Archer}{οὐκὶ πιλῆσι πρῶτά με;}{Sans me donner un baiser?}
\vers[n=1191]{ΕΥ.}{Euripide}{πάνυ γε· φίλησον αὐτόν.}{Très volontiers. \didasc{À la danseuse.} Baise-le.}
\vers[suite]{ΤΟ.}{Archer}{ὂ ὂ ὂ παπαπαπαῖ,}{Oh! oh! oh! papapapai!}
\vers[n=1192]{}{}{ὠς γλυκερὸ τὸ γλῶσσ᾽, ὤσπερ Ἀττικὸς μέλι.}{Que sa langue est douce. C'est comme du miel attique.}
\vers[n=1193]{}{}{τί οὐ κατεύδει παρ᾽ ἐμέ;}{Pourquoi ne couche-t-elle pas avec moi ?}
\vers[suite]{ΕΥ.}{Euripide}{χαῖρε τοξότα,}{Adieu, archer,}
\vers[n=1194]{}{}{οὐ γὰρ γένοιτ᾽ ἂν τοῦτο.}{cela ne se peut.}
\vers[suite]{ΤΟ.}{Archer}{ναὶ ναὶ γρᾴδιο.}{Si, si, petite vieille,}
\vers[n=1195]{}{}{ἐμοὶ κάρισο σὺ τοῦτο.}{fais-moi ce plaisir.}
\vers[suite]{ΕΥ.}{Euripide}{δώσεις οὖν δραχμήν;}{Me donneras-tu une drachme ?}
\vers[n=1196]{ΤΟ.}{Archer}{ναὶ ναῖκι δῶσι.}{Si fait, si, je te la donnerai.}
\vers[suite]{ΕΥ.}{Euripide}{τἀργύριον τοίνυν φέρε.}{Voyons ton argent.}
\vers[n=1197]{ΤΟ.}{Archer}{ἀλλ᾽ οὐκ ἔκὠδέν· ἀλλὰ τὸ συμβήνην λαβέ.}{Je n'en ai pas. Mais prends ma lance. \didasc{Il la lui confie.}}
\vers[n=1198]{}{}{ἔπειτα κομίζι σ᾽ αὐτίκ᾽. ἀκολούτι τέκνον.}{Je te l'apporte sur-le-champ. Suis-moi, mon enfant.}
\vers[n=1199]{}{}{σὺ δὲ τοῦτο τήρει τὴ γέροντο, γρᾴδιο.}{\didasc{À Euripide :} Toi, petite vieille, garde ce bonhomme.}
\vers[n=1200]{}{}{ὄνομα δέ σοι τί ἔστιν;}{Mais comment te nommes-tu ?}
\vers[suite]{ΕΥ.}{Euripide}{Ἀρτεμισία.}{Artémisia.}
\vers[n=1201]{ΤΟ.}{Archer}{μεμνῆσι τοίνυν τοὔνομ᾽· Ἀρταμουξία.}{Je le retiendrai, ce nom : Artamouxia. \didasc{Il s'en va avec la danseuse.}}
\vers[n=1202]{ΕΥ.}{Euripide}{Ἑρμῆ δόλιε ταυτὶ μὲν ἔτι καλῶς ποιεῖς.}{Hermès, dieu de la ruse, jusqu'ici tu régis tout à souhait.}
\vers[n=1203]{}{}{σὺ μὲν οὖν ἀπότρεχε παιδάριον τουτὶ λαβών·}{Toi, mon petit, enfuis-toi, avec ces instruments. \didasc{Il lui donne ses pectis.}}
\vers[n=1204]{}{}{ἐγὼ δὲ λύσω τόνδε. σὺ δ᾽ ὅπως ἀνδρικῶς}{Pendant ce temps, je le détacherai.}
\vers[n=1205]{}{}{ὅταν λυθῇς τάχιστα φεύξει καὶ τενεῖς}{\didasc{À Mnésiloque :} Quant à toi, dès que tu seras libre, vois à fuir résolument}
\vers[n=1206]{}{}{ὡς τὴν γυναῖκα καὶ τὰ παιδί᾽ οἴκαδε.}{et cours retrouver chez toi, ta femme et tes enfants.}
\vers[n=1207]{ΜΝ.}{Mnésiloque}{ἐμοὶ μελήσει ταῦτά γ᾽ ἢν ἅπαξ λυθῶ.}{J'en fais mon affaire, une fois délivré.}
\vers[n=1208]{ΕΥ.}{Euripide}{λέλυσο. σὸν ἔργον, φεῦγε πρὶν τὸν τοξότην}{Sois libre. Le reste te regarde : fuis avant que l'archer}
\vers[n=1209]{}{}{ἥκοντα καταλαβεῖν.}{ne revienne et ne t'empoigne.}
\vers[suite]{ΜΝ.}{Mnésiloque}{ἐγὼ δὴ τοῦτο δρῶ.}{C'est ce que je fais. \didasc{Euripide et Mnésiloque s'enfuient.}}
\vers[n=1210]{ΤΟ.}{Archer}{ὦ γρᾴδι᾽ ὠς καρίεντό σοι τὸ τυγάτριο,}{\didasc{Revenant en scène}. Ah! petite vieille, la charmante fillette que tu as là,}
